\chapter{Search, Filter, and Sort}

\section{Target Endpoint}

\begin{lstlisting}[language=JavaScript]
GET /api/tasks?page=1&limit=10&search=mern&status=completed&sort=createdAt
\end{lstlisting}

\begin{table}[h]
\centering
\begin{tabularx}{\textwidth}{l X}
\toprule
\textbf{Query Param} & \textbf{Purpose} \\
\midrule
\texttt{page} & Which page of results to return (used with \texttt{limit}, see Chapter~29). \\
\texttt{limit} & How many documents per page. \\
\texttt{search} & Free-text term matched against the task title (and optionally description). \\
\texttt{status} & Exact-match filter, e.g.\ \texttt{pending} / \texttt{in-progress} / \texttt{completed}. \\
\texttt{sort} & Field to sort by; a leading \texttt{-} means descending (e.g.\ \texttt{-createdAt}). \\
\bottomrule
\end{tabularx}
\end{table}

\section{Backend: Building the Filter Dynamically}

The controller builds a plain filter object incrementally, only adding keys
for parameters the client actually sent. This keeps the query flexible
without needing a different code path per combination of filters.

\begin{lstlisting}[language=JavaScript]
// backend/src/controllers/task.controller.js
const Task = require("../models/Task");

const getTasks = async (req, res, next) => {
  try {
    const { search, status, sort } = req.query;

    // ---- Build the filter object ----
    const filter = { owner: req.user.id };

    if (status) {
      filter.status = status; // exact match
    }

    if (search) {
      // Approach A: regex search (works without a text index, case-insensitive)
      filter.title = { $regex: search, $options: "i" };

      // Approach B (alternative): use a MongoDB text index instead:
      //   taskSchema.index({ title: "text", description: "text" });
      //   filter.$text = { $search: search };
      // $text is faster on large collections but requires the index above
      // and does not support partial-word matching the way $regex does.
    }

    // ---- Build the sort object ----
    // "sort=-createdAt" -> { createdAt: -1 }; "sort=title" -> { title: 1 }
    let sortObj = { createdAt: -1 }; // default: newest first
    if (sort) {
      const direction = sort.startsWith("-") ? -1 : 1;
      const field = sort.replace(/^-/, "");
      sortObj = { [field]: direction };
    }

    const tasks = await Task.find(filter).sort(sortObj);

    res.status(200).json({ data: tasks });
  } catch (error) {
    next(error);
  }
};

module.exports = { getTasks };
\end{lstlisting}

\begin{notebox}[NOTE]
This chapter combines search/filter/sort in isolation. Chapter~29 folds
\texttt{.skip()} and \texttt{.limit()} into the same query so all four
concerns --- search, filter, sort, and pagination --- work together in one
request.
\end{notebox}

\section{Frontend: Query State and Service Function}

\begin{lstlisting}[language=JavaScript]
// frontend/src/services/taskService.js
import api from "./api";

export const fetchTasks = async ({ search, status, sort }) => {
  const params = new URLSearchParams();
  if (search) params.append("search", search);
  if (status) params.append("status", status);
  if (sort) params.append("sort", sort);

  const { data } = await api.get(`/tasks?${params.toString()}`);
  return data.data;
};
\end{lstlisting}

\begin{lstlisting}[language=JavaScript]
// frontend/src/components/TaskFilters.jsx
import { useState, useEffect } from "react";
import { fetchTasks } from "../services/taskService";

function TaskFilters({ onResults }) {
  const [search, setSearch] = useState("");
  const [status, setStatus] = useState("");
  const [sort, setSort] = useState("-createdAt");

  useEffect(() => {
    const timeoutId = setTimeout(async () => {
      const tasks = await fetchTasks({ search, status, sort });
      onResults(tasks);
    }, 400); // debounce so we don't fire a request on every keystroke

    return () => clearTimeout(timeoutId);
  }, [search, status, sort, onResults]);

  return (
    <div>
      <input
        type="text"
        placeholder="Search tasks..."
        value={search}
        onChange={(e) => setSearch(e.target.value)}
      />

      <select value={status} onChange={(e) => setStatus(e.target.value)}>
        <option value="">All statuses</option>
        <option value="pending">Pending</option>
        <option value="in-progress">In Progress</option>
        <option value="completed">Completed</option>
      </select>

      <select value={sort} onChange={(e) => setSort(e.target.value)}>
        <option value="-createdAt">Newest first</option>
        <option value="createdAt">Oldest first</option>
        <option value="title">Title (A-Z)</option>
        <option value="-title">Title (Z-A)</option>
      </select>
    </div>
  );
}

export default TaskFilters;
\end{lstlisting}

\begin{tipbox}[TIP]
Debouncing the search input (waiting ~400ms after the last keystroke before
firing the request) avoids sending a network request per character typed.
\end{tipbox}
